% Includable phase-three supplement. All results here are unconditional.
\section{Lower bounds for systems of polynomial equations}
\label{p3sys:lower-bounds}

Let $K$ be a number field, and let $\mathcal O_K$ be its ring of
integers.  A polynomial system in $n$ total coordinates has the form
\[
 F_1(T,Y_1,\ldots,Y_{n-1})=\cdots
   =F_s(T,Y_1,\ldots,Y_{n-1})=0.
\]
Its degree is $\max_i\deg F_i$, including the input $T$.
Allowing several equations does not change the number of quantified
coordinates.  We also allow finitely many nonzero conditions when
stating obstructions: these specify a Zariski open subset of the common
zero set, and strengthen rather than weaken the conclusions for closed
systems.  A finite union of such presentations is treated componentwise.

The main new restriction is that \emph{two quadratic equations cannot
suffice in any number of coordinates}.  The proof includes singular,
reducible, and nonreduced intersections, rather than assuming a smooth
complete intersection.

\subsection{Integral-valued functions on simple varieties}

\begin{lemma}\label{p3sys:unirational-function}
Let $X$ be an integral $K$-variety which is $K$-unirational, let
$U\subset X$ be nonempty and open, and let $f:U\to\mathbb A^1_K$
be a regular function.  If $f(U(K))\subseteq\mathcal O_K$, then
$f$ is constant.  In particular, its image is finite.
\end{lemma}
\begin{proof}
A dominant rational parametrization reduces the assertion to a rational
function on a nonempty open subset of affine space.  If this function
were nonconstant, one of its partial derivatives would be nonzero.
Choose a $K$-point at which that derivative and every relevant
denominator are nonzero.  Such a point exists because $K$ is infinite.

Put $K_\infty=\prod_{v\mid\infty}K_v$, with one factor for each real
place and one complex factor for each pair of complex embeddings.
The diagonal copy of $K$ is dense in $K_\infty$: a $\mathbb Q$-basis
of $K$ identifies this assertion with the density of
$\mathbb Q^{[K:\mathbb Q]}$ in its real vector space.  In contrast,
$\mathcal O_K$ is a lattice, hence is closed and discrete there.
The inverse-function theorem at the chosen point, at every
archimedean place, gives a nonempty open subset of $K_\infty$ in
the local image of the rational function.  Density of $K$ makes its
$K$-values dense in this open subset.  They cannot all belong to the
discrete set $\mathcal O_K$.  This is the required contradiction.
\end{proof}

\begin{lemma}\label{p3sys:curve-function}
Let $X$ be an algebraic set over $K$ of dimension at most one,
let $U\subseteq X$ be open, and let $f:U\to\mathbb A^1_K$ be
regular.  If $f(U(K))\subseteq\mathcal O_K$, then its image is
finite.
\end{lemma}
\begin{proof}
Reduce to the finitely many irreducible components.  Zero-dimensional
components give finite sets.  If an irreducible curve is not geometrically
integral, all its $K$-points lie in the intersections of distinct
geometric components, a finite set.  Otherwise take its smooth
projective normalization $C$.  Except for finitely many points, rational
points lift to $C(K)$, and $f$ becomes a rational function on $C$.
There is nothing to prove when this function is constant.

If $C$ has genus zero and a $K$-point, it is isomorphic to
$\mathbb P^1_K$, and Lemma~\ref{p3sys:unirational-function} excludes
a nonconstant integral-valued function on its permitted open subset.
If it has no $K$-point, only the finitely many excluded points remain.
If the genus is positive, Siegel's theorem makes the set of $K$-points
where a fixed nonconstant rational function is integral finite.
For this last formulation, remove the poles of the function: the
resulting affine curve is finite over the affine line.  After inverting
a fixed finite set of primes, its coordinate ring is integral over
the ring generated by that function.  Integral function values
therefore give points integral outside this fixed set, to which
Siegel's theorem applies; see \cite[Chapter~8]{lbLang1983}.
\end{proof}

\begin{lemma}\label{p3sys:cubic-function}
Let $X\subseteq\mathbb P^N_K$ be a hypersurface defined by a
nonzero homogeneous polynomial of degree at most three.  For an open
subset $U\subseteq X$ and a regular function
$f:U\to\mathbb A^1_K$, the containment
$f(U(K))\subseteq\mathcal O_K$ implies that $f(U(K))$ is finite.
The same conclusion holds for $X=\mathbb P^N_K$.
\end{lemma}
\begin{proof}
The assertion for projective space follows from
Lemma~\ref{p3sys:unirational-function}.  Factor the defining
polynomial over $K$ and treat each irreducible factor separately;
multiplicities do not change the rational point set.  If such a
factor of degree two or three is not geometrically integral, its
geometric factors are conjugate linear forms, since these degrees
are prime.  Every $K$-point lies in their common intersection.
That intersection is a $K$-linear subspace, so the preceding lemma
for unirational varieties applies.  Linear factors are handled in
the same way.

A geometrically integral quadric with a smooth $K$-point is rational
by projection from that point.  If it has no smooth $K$-point, all
its $K$-points lie in its linear vertex.  Both cases give the result.
For a geometrically integral cubic of dimension at most one, use
Lemma~\ref{p3sys:curve-function}.  In dimension at least two,
Koll\'ar's theorem says that a cubic which is not a cone and has
a $K$-point is $K$-unirational
\cite[Theorem~2]{lbKollar2002}.  A number field is perfect,
as required in that theorem.

It remains to treat a cubic cone.  Its maximal linear vertex is
defined over $K$: in characteristic zero it is the projectivization
of the vector subspace of directions along which its defining form
is constant.  Quotient by these directions and choose a complementary
$K$-linear space.  This gives a geometrically integral cubic base
$B$ which is not a cone.  Its dimension is at least one; a
geometrically integral zero-dimensional cubic would be impossible.
Outside the vertex, projection to $B$ has affine-space fibers, whose
projective closures are the linear spaces joining a point of $B$
to the vertex.  The cone is birational to
$B\times\mathbb P^{r+1}$ if the vertex has dimension $r$.

If $B(K)$ is empty, all rational points of the cone lie in the
vertex.  If $B(K)$ is finite, each permitted fiber is an open
subset of a linear space; the function is constant on it by
Lemma~\ref{p3sys:unirational-function}.  Together with the vertex
this gives a finite image.  If $\dim B\ge2$ and $B(K)$ is
nonempty, Koll\'ar's theorem makes $B$, and hence the cone,
unirational.

The only remaining case is a cubic curve $B$ with infinitely many
$K$-points.  Pass to its normalization $C$, discarding finitely
many exceptional base points.  On a dense open of the cone we may
write $f$ as a rational function on
$C\times\mathbb P^{r+1}$.  If it depends nontrivially on the
projective-space coordinates, its specialization is nonconstant
for all base points outside a finite subset of $C$.  A rational
base point outside that subset then contradicts
Lemma~\ref{p3sys:unirational-function} on the fiber.  Thus
$f\in K(C)$.  For all but finitely many points of $C(K)$, the
permitted fiber is a nonempty open subset of projective space and
has a $K$-point.  Consequently $f$ is integral at all these base
points.  Lemma~\ref{p3sys:curve-function} forces $f$ to be
constant, because a nonconstant function on a curve has finite
fibers and $C(K)$ is infinite.  Being constant in the function
field, it is constant wherever it is regular on the irreducible
cone.  This also treats all points omitted from the birational
chart and completes the proof.
\end{proof}

\subsection{Two quadrics in any dimension}

\begin{theorem}\label{p3sys:two-quadrics-any-dimension}
Let $X\subseteq\mathbb P^N_K$ be the common zero set of at most
two homogeneous quadratic forms, let $U\subseteq X$ be open,
and let $f:U\to\mathbb A^1_K$ be regular.  If
$f(U(K))\subseteq\mathcal O_K$, then $f(U(K))$ is finite.
No smoothness, irreducibility, or reducedness assumption is made.
\end{theorem}
\begin{proof}
We induct on $N$.  The cases $N\le1$ follow from the curve lemma.
A zero equation may be discarded.  With at most one nonzero equation,
use Lemma~\ref{p3sys:cubic-function}.  If the two quadratic forms
have a common factor of degree two, they are proportional.  If their
greatest common divisor is a linear form $L$, write them as
$LL_1,LL_2$; their common zero set is the union of the linear
spaces $L=0$ and $L_1=L_2=0$.  The assertion again follows from
the unirational lemma.  We may therefore suppose that the two
forms $Q_1,Q_2$ are coprime.

If $U(K)$ is empty there is nothing to prove.  Choose $p\in U(K)$
and coordinates with $p=(1:0:\cdots:0)$.  Since $Q_i(p)=0$,
we can write
\[
 Q_1=X_0L_1(X_1,\ldots,X_N)+q_1(X_1,\ldots,X_N),\qquad
 Q_2=X_0L_2(X_1,\ldots,X_N)+q_2(X_1,\ldots,X_N),
\]
where the $L_i$ are linear and the $q_i$ are quadratic.  We consider
the rank of the span of $L_1,L_2$.

If the rank is two, projection from $p$ is an isomorphism between
\[
 X\setminus V(L_1,L_2)
 \quad\text{and}\quad
 V(L_1q_2-L_2q_1)\setminus V(L_1,L_2)
          \ \subseteq\mathbb P^{N-1}_K.
\]
On the chart $L_i\ne0$ the inverse is given by
$X_0=-q_i/L_i$; these formulas agree on the overlap.  The cubic
$L_1q_2-L_2q_1$ is nonzero.  Otherwise coprimality of the two
linear forms would give $q_1=L_1M$, $q_2=L_2M$ for a linear
form $M$, making $X_0+M$ a common factor of $Q_1,Q_2$.
The image of $f$ on this open part is finite by
Lemma~\ref{p3sys:cubic-function}.  The exceptional part
$X\cap V(L_1,L_2)$ is cut out by two quadrics in the smaller
linear space $V(L_1,L_2)\simeq\mathbb P^{N-2}_K$.
The induction hypothesis applies to the restriction of $f$ and
the open subset inherited from $U$.

If the rank is one, replace the two forms by an invertible linear
combination so that $L_1\ne0$ and $L_2=0$.  Away from $L_1=0$,
projection from $p$ is an isomorphism to the open subset
\[
                         V(q_2)\setminus V(L_1)
                          \subseteq\mathbb P^{N-1}_K,
\]
with inverse $X_0=-q_1/L_1$.  The polynomial $q_2$ is nonzero,
since the second equation was nonzero.  The quadric lemma gives
a finite image here.  On the exceptional linear space $L_1=0$
we again have two quadrics in a smaller projective space, so
induction gives a finite image there as well.

If the rank is zero, both equations are independent of $X_0$,
and $p$ is a vertex of the intersection.  For every $q\in U(K)$,
the line joining $p$ to $q$ lies in $X$.  The restriction of $f$
to the open part of this line has only integral values at its
$K$-points, hence is constant by the unirational lemma.  It is
defined at both $p$ and $q$, so $f(q)=f(p)$.  The image is a
singleton.  The three cases prove the induction.
\end{proof}

\begin{corollary}\label{p3sys:two-quadratics-affine}
A system of at most two equations of total degree at most two,
in any number of rational coordinates over $K$, cannot define
$\mathcal O_K$ in its input.  More generally, whenever its
projection, possibly restricted by finitely many polynomial
nonzero conditions, lies in $\mathcal O_K$, that projection is
finite.  The same holds for a finite union of such presentations.
\end{corollary}
\begin{proof}
Homogenize each equation to degree two and take the open chart
where the homogenizing coordinate is nonzero.  The input becomes
a regular rational function on that chart.  Apply the theorem;
a finite union of finite images is finite.
\end{proof}

Linear equations can first be solved by Gaussian elimination.  On the
resulting affine subspace, the input is an affine linear function and
the remaining residuals are still quadratic.  Thus the same obstruction
applies whenever those residuals span a vector space of dimension at
most two.  A quadratic system defining $\mathcal O_K$ must have at
least three independent quadratic residuals after the linear equations
have been removed.  This assertion counts equations, not variables.

\subsection{Restrictions on the number of coordinates}

\begin{lemma}\label{p3sys:horizontal-closure}
Let $U$ be a locally closed subset of affine space over $K$ and
let $f:U\to\mathbb A^1_K$ be regular.  If $f(U(K))$ is an
infinite subset of $\mathcal O_K$, the Zariski closure of $U(K)$
has an irreducible component $Z$ of dimension at least two,
with Zariski-dense $K$-points and nonconstant input function.
The component is geometrically integral.
\end{lemma}
\begin{proof}
Take the finitely many irreducible components of the closure.  The
rational points are dense on each component after removing its
intersections with the others.  A component that is not geometrically
integral has all its $K$-points in the intersections of distinct
geometric components; this contradicts density in characteristic
zero.  If every component had dimension at most one or constant
input, the curve lemma would make the union of their input images
finite.  At least one component has the stated properties.
\end{proof}

\begin{theorem}\label{p3sys:small-coordinate-bounds}
Suppose a finite polynomial system, possibly with nonzero guards,
has an infinite rational input projection contained in
$\mathcal O_K$.
\begin{enumerate}
\item It uses at least three total coordinates.  This holds at every
degree and for a finite union of presentations.
\item If it uses three total coordinates, every nonzero residual
has degree at least four.  In fact the residuals have a common
irreducible factor of degree at least four with nonconstant input.
\item If it uses at most four total coordinates, its residuals
cannot all be quadratic.  More strongly, in four coordinates it
cannot contain two coprime residuals of degree at most two,
regardless of the degrees of its other residuals.
\end{enumerate}
\end{theorem}
\begin{proof}
For the first assertion, a proper algebraic subset of affine
two-space has dimension at most one, so the curve lemma applies.
If the closure is all of affine two-space, a locally closed subset
with that closure contains a nonempty open subset; its input cannot
be integral-valued there by the unirational lemma.  The cases with
fewer coordinates are included.  A finite union does not help,
since an infinite image would require an infinite image from one
member.

In three coordinates, choose the component $Z$ in
Lemma~\ref{p3sys:horizontal-closure}.  If there is a nonzero
residual, its common zero set is proper, so $\dim Z=2$.  Its
height-one prime ideal is generated by an irreducible polynomial
$g$.  Every residual vanishes on $Z$, hence is divisible by $g$.
The permitted open subset of $V(g)$ is contained in the source
system and has dense rational points and nonconstant input.
Lemma~\ref{p3sys:cubic-function} excludes $\deg g\le3$.
This proves the second assertion.  If all residuals are zero,
the source is open in affine space and is excluded by the
unirational lemma.

For the stronger statement in four coordinates, suppose two
residuals $Q_1,Q_2$ of degree at most two are coprime.  Their
common zero set has dimension at most two.  Thus the horizontal
component $Z$ again has dimension two and is an irreducible
component of this common zero set.  Homogenize to quadrics,
restrict to the affine chart, and remove all other components.
Intersect further with the source's nonzero guards.  This gives
an open subset of the intersection of two projective quadrics
whose rational points lie in the source system.  The image is
finite by Theorem~\ref{p3sys:two-quadrics-any-dimension}, whereas
density of the points on $Z$ and nonconstancy of the input make
that image infinite.  This is a contradiction.

Finally suppose all residuals have degree at most two.  A nonzero
constant makes the source empty.  A nonzero linear equation reduces
the ambient affine dimension to at most three and is excluded by
the second assertion.  If the quadratic residuals have a common
factor, either they are proportional quadrics or they are a common
linear factor times linear forms.  Their common zero set is then
a quadric or a union of linear spaces, and the cubic-function
lemma applies on every permitted open part.  Otherwise choose
one residual $Q_1$ and a linear combination $Q_2$ of the others
which avoids every irreducible factor of $Q_1$.  Such a choice
exists over the infinite field $K$: each forbidden divisibility
condition is a proper linear subspace of the coefficient choices.
The resulting pair is coprime, so the preceding paragraph applies.
\end{proof}

In particular, a quadratic system defining the integers over the
rationals needs at least five total coordinates and at least three
independent quadratic equations after removal of linear constraints.
A finite system of cubic equations in three total coordinates cannot
define the integers either.  These restrictions leave open systems in
the requested range of five to ten coordinates.  They do not provide
a construction in that range and do not exclude it.
